% =====================================================================
\section{Read Me First}
\label{sec:readme}
% =====================================================================

\subsection{The project in one paragraph}

The base paper (Golestani et al., IEEE Access 2022) proposes a PD-type attitude controller that keeps \emph{both} the
attitude quaternion and the angular velocity of a rigid body inside time-varying bounds that close in a finite, chosen
time, while the actuators saturate and an unknown disturbance acts. Its selling point is simplicity: an ordinary PD law
plus one barrier-weighted damping term and one scalar adaptive estimate. We (1) implemented every equation of the
paper, (2) reproduced both of its simulation studies quantitatively, (3) wrote an independent MATLAB toy model that
re-derives every block and re-implements the experiments, (4) tested the controller beyond the paper (300-run Monte-Carlo studies, parameter
sensitivity, an alternative barrier form), and (5) integrated it on the Gazebo X3 quadcopter through ROS~2, where we
found and fixed the faults that made our first version crash.

\subsection{Results at a glance}

\begin{center}\small
\begin{tabularx}{\linewidth}{@{}>{\raggedright\arraybackslash}p{4.6cm} X >{\raggedright\arraybackslash}p{2.2cm}@{}}
\toprule
\textbf{Question} & \textbf{Answer (measured)} & \textbf{Where}\\
\midrule
Does our implementation reproduce the paper's Part~1? & Yes. For $T_s=15$\,s and $10$\,s both $\qv$ and $\vw$ stay strictly inside their
envelopes (0\,\% of the time outside; peak $|\varepsilon|=0.871/0.912$ and $0.924/0.942$), and the shorter $T_s$ converges faster
(settling 5.25\,s vs.\ 4.53\,s), as the paper states. & \cref{sec:part1}\\
Does it reproduce Part~2 (constrained vs.\ conventional PD)? & Yes. Same gains, limit, disturbance and initial state: the constrained law
never leaves the envelopes; conventional PD leaves the angular-velocity envelope 27.5\,\% of the time. The constrained
law also uses 36\,\% less torque energy and halves the final rate error. & \cref{sec:part2}\\
Is this a one-off? & No. 300 random initial states in the paper's regime: 97.0\,\% violation-free (95\,\% CI 94.4--98.4) for
the constrained law, 0.0\,\% for PD. & \cref{sec:mc}\\
Can it be checked independently? & The MATLAB toy model re-implements the spacecraft experiments separately
(\code{toy\_model/step6}); \code{docs/make\_crosscheck.py} puts its numbers next to the Python ones (\cref{tab:crosscheck}). & \cref{sec:crosscheck}\\
What did we find that the paper does not say? & Four unstated parameters; Part~2 starts on its own boundary; the barrier is
one-sided; the envelope must be feasible for the actuator; Theorem~1 proves convergence, not envelope invariance. &
\cref{sec:findings}\\
Does it fly a quadcopter? & Yes, after fixing our first node. The paper's law, re-sized for the X3, hovers for 180\,s and flies
a $2\times2$\,m square (4.0\,cm RMS error) on a software-in-the-loop model built from Gazebo's X3 files, with 0\,\% attitude-envelope
violations. The ROS~2 node \code{gazebo.py} runs exactly this controller. & \cref{sec:x3res}\\
Why did the first Gazebo version crash? & Two faults, isolated one by one: a mixer whose signs did not match the X3
(no pitch authority) and an altitude-rate estimate that was never updated. Fixing both gives a stable flight. & \cref{sec:x3fault}\\
\bottomrule
\end{tabularx}
\end{center}

\subsection{The documents and how to read them}

\begin{center}\small
\begin{tabularx}{\linewidth}{@{}l X@{}}
\toprule
\textbf{Document} & \textbf{Content}\\
\midrule
\code{docs/complete\_documentation/} (this PDF) & project, results, code walk-through, X3 integration, claims\\
\code{docs/mathematical\_models/} & every equation: boxed statement, symbol table, derivation, hand-checkable toy example, flow diagrams\\
\code{toy\_model/} & MATLAB scripts (\code{step1}--\code{step8}) that print the toy examples and plot each idea; \code{run\_all.m} runs all of them\\
\code{docs/results/*.md, *.json} & every number of this document, written by the scripts\\
\code{docs/figures/} & every figure of this document, written by the scripts\\
\bottomrule
\end{tabularx}
\end{center}

\textbf{Colour code of the boxes.} \textcolor{cX3}{Blue}: the idea in plain words. \textcolor{cEnv}{Orange}: a toy
example with numbers. \textcolor{cOk}{Green}: what it means in our project. \textcolor{cWarn}{Red}: something to be
careful about.

\textbf{Naming of the controllers.} Throughout, \emph{constrained (paper Eq.~13)} is the paper's law exactly as printed;
\emph{constrained (symmetric barrier)} is our variant with $(1-\varepsilon^2)^{-2}$ in place of $(1-\varepsilon)^{-2}$
(\cref{sec:barrierform}); \emph{conventional PD} is the same law with $\vu_{cons}=0$, i.e.\ the paper's Part~2 baseline.
